\begin{frame}[fragile]{Why not 1,2,3 but 4,4,4?}

\begin{itemize}
	\item When we use the setTimeout() function, we should be careful when we use it with the var modifier.
    \item When we execute this code, the output is not what we expect (1, 2, 3), but (4, 4, 4). 
    \item This is confusing, especially for those who are familiar with multi-thread programming languages. 
\end{itemize}

\begin{columns}%
\begin{column}{0.65\textwidth}%
\begin{Code}{}%
\begin{lstlisting}[firstnumber=1, label=glabels, xleftmargin=0pt, basicstyle=\tiny]% 

for (var i = 1; i <= 3; i++) {
  setTimeout(function(){ console.log(i); }, 0);
};
\end{lstlisting}%   
\end{Code}%
\end{column}%
\begin{column}{0.3\textwidth}%
\begin{center}%
\begin{Command}{Results}%
\begin{lstlisting}[firstnumber=1, label=glabels, xleftmargin=0pt, basicstyle=\tiny]% 

4
4
4
\end{lstlisting}%   
\end{Command}%
\end{center}%
\end{column}%
\end{columns}%
\end{frame}

\begin{comment}
\begin{frame}[fragile]{}

\begin{itemize}
	\item When we use `chatGPT' to find a bug in this code, it points out that the error is from using `var' and how the callback queue (message queue) works. 
\end{itemize}

\begin{columns}%
\begin{column}{0.6\textwidth}%
\textit{In the code, all setTimeout functions are added to the queue almost immediately, and they are all set to run after a delay of 0 milliseconds. However, the for loop will complete before the first setTimeout function is executed, so by the time the functions are executed, the value of i will be 4 in all cases.}
\end{column}%
\begin{column}{0.3\textwidth}%
\begin{center}%
\begin{Code}{}%
\begin{lstlisting}[firstnumber=1, label=glabels, xleftmargin=0pt, basicstyle=\tiny]% 

for (var i = 1; i <= 3; i++) {
  (function(j) {
    setTimeout(
      function(){ 
        console.log(j); 
      }, 
      0);
  })(i);
};
\end{lstlisting}%   
\end{Code}%
\end{center}%
\end{column}%
\end{columns}%
\end{frame}
\end{comment}
\begin{frame}[fragile]{}

\begin{itemize}
	\item It is because `JavaScript event handlers don't run until the thread is free.`
	\item In other words, because of its single-thread nature, the functions in the callback queue cannot be executed until the for loop is finished.
	\item And because of the nature of the var, the variable `i' (with a value of 4) makes all the `console.log(i)' code print out 4. 
	\item This is a tricky bug that is hard to understand unless we are familiar with the bad parts of JavaScript. 
\end{itemize}
\end{frame}

\begin{frame}[fragile]{Another Example}

\begin{itemize}
	\item In this code, if JavaScript uses the multi-thread model, The console.log() function will be executed in 500 ms.
	\item However, with a single thread model, JavaScript while loop (line 6) cannot be interrupted for 1000ms; as a result, the console.log() is executed after 1000ms. 
\end{itemize}

\begin{Code}{}%
\begin{lstlisting}[firstnumber=1, label=glabels, xleftmargin=0pt, basicstyle=\tiny]% 

var start = new Date;
setTimeout(function(){
  var end = new Date;
  console.log('Time elapsed:', end - start, 'ms');
}, 500);
while (new Date - start < 1000) {};
\end{lstlisting}%   
\end{Code}%
\begin{Command}{Results}%
\begin{lstlisting}[firstnumber=1, label=glabels, xleftmargin=0pt, basicstyle=\tiny]% 

Time elapsed: 1002ms
\end{lstlisting}%   
\end{Command}%

\end{frame}

\begin{frame}[fragile]{The Inaccuracies of setTimeOut}

\begin{itemize}
	\item We use the setInterval() function with 0 argument to execute the lambda expression given as the first argument as quickly as possible. 
\end{itemize}

\begin{Code}{}%
\begin{lstlisting}[firstnumber=1, label=glabels, xleftmargin=0pt, basicstyle=\scriptsize]% 

var fireCount = 0;
var start = new Date;
var timer = setInterval(function() {
  if (new Date - start > 1000) {
    clearInterval(timer);
    console.log(fireCount);
    return;
  }
  fireCount++;
}, 0); // 843
\end{lstlisting}%   
\end{Code}%
\end{frame}

\begin{frame}[fragile]{}

\begin{itemize}
	\item The setInterval() is the same as the setTimeOut() function as it executes the given lambda expression after the specified time (the argument given as the second argument). 
	\item The printed output is about 843; considering we expect about 4.5 million with a loop, this is shockingly slow.  
\end{itemize}

\begin{Code}{}%
\begin{lstlisting}[firstnumber=1, label=glabels, xleftmargin=0pt, basicstyle=\scriptsize]% 

function test() {
  while (true) {
    if (new Date - start > 1000) {
      console.log(fireCount);
      return;
    }
    fireCount++;
  }
} 
test(); // 4522448
\end{lstlisting}%   
\end{Code}%

\end{frame}